\section{Methodology}
\label{sec:methodology}

\subsection{Portfolio construction and notation}
Consider $n$ risky assets. Let $r_t \in \mathbb{R}^n$ denote the vector of simple returns on trading day $t$, and let $w \in \mathbb{R}^n_{\ge 0}$ denote portfolio weights satisfying $\sum_i w_i = 1$ after normalisation. The portfolio return is
\begin{equation*}
r_{p,t} = w^\top r_t .
\end{equation*}
Vectors $w$ are supplied exclusively from the validated interface state and are not inferred by the language model, thereby avoiding implicit reweighting such as undocumented equal-weight assumptions.

\subsection{Risk and performance metrics (current build)}
The implemented pipeline computes the following quantities on the sample of portfolio and asset returns:
\begin{itemize}
    \item \textbf{Annualised volatility.} The sample standard deviation of $\{r_{p,t}\}$, multiplied by $\sqrt{252}$ under the conventional mapping from daily to annual units for equities.
    \item \textbf{Historical quantile loss.} The empirical 5th percentile of $\{r_{p,t}\}$ at a one-day horizon, reported without scaling to multi-day regulatory value-at-risk horizons; it constitutes a descriptive tail measure rather than a full loss-distribution model.
    \item \textbf{Sharpe-type ratio.} The ratio of mean daily excess return to standard deviation of daily returns, annualised by $\sqrt{252}$, with the risk-free rate set identically to zero for tractability; interpretability as a true economic Sharpe ratio therefore requires mental adjustment for the user's financing rate.
    \item \textbf{Maximum drawdown.} Computed on the compounded wealth process $\prod_s (1+r_{p,s})$ as the largest relative decline from a running maximum to a subsequent trough.
    \item \textbf{Average pairwise correlation.} The arithmetic mean of off-diagonal entries in the sample correlation matrix of asset simple returns.
    \item \textbf{Concentration (Herfindahl).} $\sum_i w_i^2$, equal to unity under perfect concentration and to $1/n$ under equal weighting.
\end{itemize}
The six scalars enter a piecewise rule-based mapping to a composite score in $[0,1]$ and to ordered risk categories (\texttt{Low}, \texttt{Medium}, \texttt{High}). Thresholds are chosen for demonstration and narrative coherence; they are not calibrated against regulatory templates.

The intent classifier vocabulary may reference additional statistics (e.g.\ Sortino ratio, beta, conditional tail expectations) that constitute planned extensions; any metric not produced by the executed code path should be explicitly identified as absent in empirical sections of this report.

\subsection{Prospective machine-learning components}
The software repository permits future integration of sequence models (e.g.\ recurrent networks for forward-looking volatility) and learned risk taxonomies. At the time of writing, categorical risk bands rely on deterministic rules rather than trained classifiers. Adoption of learning-based modules would oblige a complementary discussion of training and validation protocols, information leakage across time, and the provisional nature of forecast outputs.

\subsection{Retrieval-augmented generation}
The RAG subsystem partitions background knowledge into thematic stores (including, inter alia, ticker disambiguation aids, firm-level summaries, macroeconomic indicators, curated definitions, and strategic-allocation notes). Retrieval is conditioned on the inferred dialogue intent: only subsets of corpora are queried for a given turn. Candidate passages undergo deduplication and scoring; retrieval itself employs a hybrid of lexical matching (Okapi BM25) and dense retrieval of sentence-transformer embeddings persisted in Chroma. Logging supports retrospective measurement of retrieval quality (e.g.\ recall at fixed $k$).

In the \textbf{reference Streamlit deployment} (\texttt{chatbot.py}), assistant turns are produced solely by deterministic Python: template Markdown, numerically computed metrics, and optional retrieved excerpts concatenated in the orchestrator. There is \emph{no} second large-language-model call that rewrites the full answer in that path; richer ``answer models'' may exist in ancillary scripts or future revisions, but they are not required for the live loop described here. Relative to unconstrained generation, retrieved spans narrow---without eliminating---hallucination risk by anchoring visible text to locally stored evidence, subject to corpus coverage and chunk boundaries.

\subsection{Agent and intent orchestration}
Each utterance receives a primary intent label from Gemini among holistic portfolio analysis, single-metric interrogation, conceptual explanation, trend or outlook questioning, conversational follow-up, and general dialogue. Optional secondary intents and normalised entity fields (metrics, concepts) may accompany the classification.

Dispatch logic implemented in Python branches on the primary label. Computationally intensive paths reuse a single price fetch and return construction step where feasible; explanatory paths may issue retrieval calls with intent tags recognised by the RAG layer (\texttt{full\_analysis}, \texttt{concept\_explanation}, \texttt{trend\_prediction}). Follow-up handling surfaces the preceding assistant message from session history rather than maintaining long-horizon episodic memory. Separation of probabilistic classification from deterministic numerical execution is intentional: portfolio arithmetic remains auditable code, not prompt-dependent improvisation.
